The third generation iPod Nano uses a DFU installer, \fname{mks5lboot},
to install the bootloader in flash. Do not use \fname{ipodpatcher}.

\warn{Rockbox supports the Nano 3G models fitted with Hynix NAND, the
4~GB A514D3AD and the 8~GB A555D5AD. Apple fitted several other NAND
chips to this player. Rockbox drives only chips that have been tested on
hardware: on any other chip the bootloader says the NAND is not supported
and starts Apple's firmware instead, leaving the player's storage
untouched. The port is experimental: Apple's proactive
refresh of blocks needing ECC correction is not implemented, and
bad-block remapping has not been exercised on hardware. Keep a backup of
all media.}

If your player has another chip, you can help have it added. The NAND check in \fname{utils/ipodnano3g/nandcheck} of the
source tree runs from DFU mode, installs nothing, writes nothing to the
player, and collects what validating the chip needs: see that
directory's \fname{README.md}.

\caps{Rockbox Utility} can install the bootloader. To install it by hand
instead, follow the steps below.

Use \fname{mks5lboot} built from Rockbox sources
with Nano 3G support, and the Nano 3G bootloader
\fname{bootloader-ipodnano3g.ipod}. Older iPod Classic installers do not
support this player. Build instructions and USB driver requirements are
in \fname{utils/mks5lboot/README} in the source tree.

\begin{enumerate}
\item Install the Rockbox files on the player's FAT32 volume before
installing the bootloader.
\item Connect the player by USB. Close iTunes and stop iTunesHelper if
they are running.
\item Run \fname{mks5lboot --dfuscan --loop} in a terminal, then put the
player in
\href{https://freemyipod.org/wiki/Modes#DFU_mode}{DFU mode}.
When it is detected, stop the scan with Ctrl+C.
\item Run \fname{mks5lboot --bl-inst path/to/bootloader-ipodnano3g.ipod},
replacing the path with the location of your Nano 3G bootloader.
\item Wait for the installer to finish and the player to restart.
\end{enumerate}
